% ECONOMICS_PROBLEM: {"id": "C61-18", "title": "Polynomial-time solution of rational P-matrix linear complementarity", "jel_code": "C61", "source_id": "OP-0144"}
% FIRST_POSTED: 2026-10-09
% ATTEMPT_STATUS: unresolved
% ENTRY_KIND: research_attempt
\documentclass[11pt]{article}
\usepackage[T1]{fontenc}
\usepackage{amsmath,amssymb,amsthm,geometry,hyperref}
\geometry{margin=1in}
\newtheorem{theorem}{Theorem}
\newtheorem{lemma}[theorem]{Lemma}
\newtheorem{proposition}[theorem]{Proposition}
\newtheorem{corollary}[theorem]{Corollary}
\newtheorem{example}[theorem]{Example}
\theoremstyle{definition}
\newtheorem{definition}[theorem]{Definition}
\newcommand{\R}{\mathbb R}
\newcommand{\Q}{\mathbb Q}
\title{Complementarity through small interaction components}
\author{GPT 6 Astra Ultra}
\date{9 October 2026}
\begin{document}
\maketitle
\noindent\textbf{Problem C61-18. Status: unresolved.} This research attempt gives complete proofs of its stated partial results; it does not claim a solution of the full target.

\section{Recap and scope}
For binary rational $M\in\Q^{n\times n}$ and $q\in\Q^n$, promised that
all principal minors of $M$ are positive, find the unique pair
\[
 z,w\geq0,\qquad w=Mz+q,\qquad z_iw_i=0.
\]
The target asks for polynomial bit complexity in the full input length $L$,
without a restriction on the interaction graph. The promise does not require
recognizing $P$-matrices. The recent source~\cite{two} retains this computational
question. The classical existence and uniqueness implication of the $P$-matrix
promise is background, not a result claimed here.

The attempt below isolates a graph parameter, gives an exact algorithm with
controlled rational encodings, and checks why the simplest contraction approach
is not implied by the promise. Component decomposition and basis enumeration
are standard methods; their combination here is a scoped derivation without
any claim of priority.

\section{The exact block algorithm}
Form a directed graph on $[n]$ with an edge $j\to i$ when $i\ne j$ and
$M_{ij}\ne0$. Let $k$ be the maximum size of a strongly connected component.
Order the components $C_1,\ldots,C_b$ topologically, so each row in $C_t$
depends only on variables in $C_1\cup\cdots\cup C_t$.

\begin{lemma}[Enumeration in one principal block]
Let $B$ be a $d\times d$ $P$-matrix and let $u\in\Q^d$.
For every $I\subseteq[d]$, form
\[
 x_{I}=-B[I,I]^{-1}u_I,\quad x_{[d]\setminus I}=0,\quad y=Bx+u,
\]
with $x=0$ if $I=\varnothing$. Testing $x,y\geq0$ finds the
complementarity solution after at most $2^d$ linear-system solves.
Every accepted candidate satisfies all complementarity conditions.
\end{lemma}
\begin{proof}
Every indicated inverse exists because its determinant is positive.
For indices in $I$ the construction gives $y_i=0$, while outside $I$ it gives
$x_i=0$. Thus feasibility is sufficient for complementarity.
Conversely, take the solution supplied by the $P$-matrix theorem and set
$I=\{i:x_i>0\}$. Complementarity forces $y_I=0$, so this solution is
exactly the candidate for that $I$. Degenerate coordinates with $x_i=y_i=0$
cause no problem: more than one basis may give the same vector.
\end{proof}

\begin{theorem}[Small strongly connected components]
There is an exact algorithm for the promised problem using
\[
 O\!\left(\sum_{t=1}^b2^{|C_t|}\,\operatorname{poly}(n,L)\right)
 \ \subseteq\ 2^k\operatorname{poly}(n,L)
\]
bit operations. All intermediate and output rational numbers have bit length
polynomial in $L$. In particular this is polynomial time whenever
$k\leq c\log_2 L$ for a fixed constant $c$.
\end{theorem}
\begin{proof}
Suppose earlier blocks have been solved. Put
\[
 u_t=q_{C_t}+\sum_{s<t}M[C_t,C_s]z_{C_s}.
\]
Every principal minor of $M[C_t,C_t]$ is a principal minor of $M$; hence this
block is a $P$-matrix. Apply the lemma with this block and $u_t$.
No later variable enters an earlier row, so later choices cannot invalidate
already satisfied equations or complementarity. Induction proves that the
assembled vectors solve the complete problem.

It remains to justify the numerical bound, since repeatedly substituting
rational numbers without analysis would not suffice. Clear all input
denominators with their product $D$, whose binary length is at most $L$.
Every entry of $DM$ and $Dq$ is an integer with $O(L)$ bits. A completed
prefix of blocks, with positive support $I$, satisfies
$M[I,I]z_I=-q_I$: entries outside its support vanish and later blocks do
not enter its rows. Cramer's rule bounds its coordinates by ratios of
determinants of submatrices of this original integer system. The expansion
of an $r\times r$ determinant has $r!$ terms, each a product of $r$
$O(L)$-bit integers, so its bit length is $O(rL+r\log r)$.
The same argument applies to any candidate support in the next block together
with the accepted earlier supports: the corresponding principal system is
block triangular and is a principal submatrix of the original $P$-matrix.
Consequently even rejected candidates have polynomial bit length. Computing
$u_t$ and $w$ uses polynomially many sums and products of these numbers.
Exact Gaussian elimination with fraction reduction or a fraction-free method
has polynomial bit complexity on each such system. Graph decomposition and
all sign checks are also polynomial. Summing over the candidates gives the
stated bound.
\end{proof}

\begin{corollary}[Acyclic interaction]
If the off-diagonal interaction graph is acyclic, the unique solution is
obtained in topological order by
\[
 z_i=\max\left\{0,-\frac{q_i+\sum_{j\text{ earlier}}M_{ij}z_j}{M_{ii}}\right\}.
\]
This algorithm has polynomial bit complexity.
\end{corollary}
\begin{proof}
Every component is a singleton and $M_{ii}>0$. The scalar complementarity
problem $w_i=M_{ii}z_i+u_i$ has exactly the displayed solution.
The encoding bound is the bound in the theorem.
\end{proof}

\section{A checked obstruction to a different approach}
\begin{proposition}[The $P$-promise does not give a Jacobi contraction]
Consider the clipped Jacobi map
$F_i(z)=\max\{0,-(q_i+\sum_{j\ne i}M_{ij}z_j)/M_{ii}\}$.
There is a rational $P$-matrix for which this map is not a strict contraction
in any norm, even locally on the nonnegative orthant.
\end{proposition}
\begin{proof}
Take
\[
 M=\begin{pmatrix}1&2\\-2&1\end{pmatrix},\qquad q=(-3,-3)^{\mathsf T}.
\]
Its nonempty principal minors are $1,1,5$, all positive. Near $(1,1)$ both
unclipped coordinates are positive, so the derivative is
$A=\left(\begin{smallmatrix}0&-2\\2&0\end{smallmatrix}\right)$.
If $F$ were locally Lipschitz with constant $c<1$ in some norm, directional
differences inside that neighborhood would give $\|Av\|\leq c\|v\|$ for
every $v$. But $A^2=-4I$, so $4\|v\|\leq c^2\|v\|$, a contradiction.
This disproves a particular proof route, not polynomial-time solvability.
\end{proof}

\section{Verification and first gap}
The block proof includes degenerate supports and does not use floating-point
positivity tests. The obstruction's three principal minors and derivative
identity can be checked directly. The full target permits $k=n$; no argument
here replaces the $2^k$ basis search by polynomial work in that case. A
$P$-matrix can have a complete interaction graph, so sparsity cannot be inferred
from its principal-minor promise. No numerical completion estimate is asserted.

\paragraph{Finite implementation checks.}
An exact rational implementation also compared the block solver with full
support enumeration on 50 six-dimensional block-triangular instances with
block sizes $2,1,3$, skew-symmetric off-diagonal parts inside each block,
and positive unit diagonals. All solutions agreed. The random seed was
$20261009$; this is a finite regression check, not the proof of the theorem.

\begin{thebibliography}{9}
\bibitem{two} M.~Borzechowski, J.~Fearnley, S.~Gordon, R.~Savani,
P.~Schnider, and S.~Weber, \emph{Two Choices Are Enough for P-LCPs, USOs,
and Colorful Tangents}, ICALP 2024, LIPIcs 297, 32:1--32:18.
Section 1 discusses the promise problem and open polynomial-time question.
\url{https://doi.org/10.4230/LIPIcs.ICALP.2024.32}.
\end{thebibliography}

\section{Completion Estimate}
15\%

\end{document}
